
\documentclass[12pt,a4paper]{article}

\usepackage[a4paper,margin=1in]{geometry}
\usepackage{graphicx}
\usepackage{booktabs}
\usepackage{amsmath,amssymb}
\usepackage{float}
\usepackage{hyperref}
\usepackage{xcolor}
\usepackage{listings}
\usepackage{enumitem}
\usepackage{fancyhdr}
\usepackage{longtable}

\hypersetup{colorlinks=true,linkcolor=blue,urlcolor=blue}

\pagestyle{fancy}
\fancyhf{}
\lhead{ICS1512}
\rhead{Machine Learning Algorithms Laboratory}
\cfoot{\thepage}

\lstset{
language=Python,
basicstyle=\ttfamily\small,
numbers=left,
frame=single,
breaklines=true,
showstringspaces=false
}

\begin{document}

\begin{tabular}{|p{4cm}|p{11cm}|}
\hline
Experiment No. & 6 \\ \hline
Experiment Title & Bagging, Boosting, and Stacked Ensemble Models \\ \hline

Student Name & Naren Karthik Kandasamy \\ \hline
Register Number & 3122247001036 \\ \hline
Faculty & Dr. Poreddy Ajay Kumar Reddy \\ \hline
Submission Date & 31/8/2026 \\ \hline
\end{tabular}

\section{Objective}
To understand and implement ensemble learning strategies including Bagging, Boosting, and Stacking. To compare these models in terms of accuracy, stability, and generalization, and analyze their effect on bias and variance.

\section{Problem Statement}
The problem involves classifying tumor samples from the Wisconsin Diagnostic Breast Cancer dataset. We employ heterogeneous ensemble models to combine multiple base learners and theoretically compare how Bagging reduces variance while Boosting addresses model bias.

\section{Dataset Description}
\begin{longtable}{|p{6cm}|p{8cm}|}
\hline
Dataset Name & Wisconsin Diagnostic Breast Cancer \\ \hline
Dataset Source & UCI Machine Learning Repository \\ \hline
Number of Samples & 569 \\ \hline
Number of Features & 30 \\ \hline
Number of Classes & 2 (Malignant, Benign) \\ \hline
Missing Values & None \\ \hline
Train-Test Split & 80-20 \\ \hline
\end{longtable}

\section{Exploratory Data Analysis}
Include all relevant visualizations.
\begin{itemize}
\item Dataset overview
\item Statistical summary
\item Missing value analysis
\item Class distribution
\item Correlation matrix
\item Feature distribution
\end{itemize}

\subsection*{Figure Template}



\textbf{Inference:}

Exploratory Data Analysis was thoroughly conducted in the previous experiment. The data distributions and correlations remain identical, so we proceed directly to modeling.

\section{Data Preprocessing}

Explain:
\begin{itemize}
\item Missing value handling
\item Encoding
\item Feature scaling
\item Train-test split
\item Feature engineering (if applicable)
\end{itemize}

\section{Machine Learning Algorithm}
Explain:
\begin{itemize}
\item Working principle
\item Mathematical formulation
\item Advantages
\item Limitations
\end{itemize}

\section{Experimental Setup}
\begin{tabular}{ll}
\toprule
Python Version & \\
Scikit-Learn Version & \\
IDE & \\
Operating System & \\
Hardware & \\
\bottomrule
\end{tabular}

\section{Hyperparameter Selection}

\begin{tabular}{ll}
\toprule
Parameter & Value\\
\midrule
 & \\
 & \\
\bottomrule
\end{tabular}



\section{Results}

\begin{tabular}{lc}
\toprule
Metric & Value\\
\midrule
\textbf{Table 4: Performance Comparison of Ensemble Models}\\\\ \begin{tabular}{lllll}\toprule Model & Accuracy (\%) & Precision & Recall & F1 Score\\ \midrule Bagging & 96.49 & 0.9589 & 0.9859 & 0.9722\\ AdaBoost & 96.49 & 0.9589 & 0.9859 & 0.9722\\ Gradient Boosting & 95.61 & 0.9583 & 0.9718 & 0.9650\\ Stacked Ensemble & 96.49 & 0.9589 & 0.9859 & 0.9722\\ \bottomrule \end{tabular}
\bottomrule
\end{tabular}

\section{Result Visualization}

Include all graphs generated during the experiment.

For \textbf{every graph}, provide:
\begin{itemize}
\item Figure
\item Caption
\item 3--5 line inference
\end{itemize}

Recommended figures:
\begin{enumerate}
\item Correlation Matrix
\item Class Distribution
\item Confusion Matrix
\item ROC Curve
\item Precision-Recall Curve
\item Learning Curve (if applicable)
\item Feature Importance (if applicable)
\end{enumerate}
\begin{figure}[H]
\centering
\includegraphics[width=0.8\linewidth]{Ex6/ROC_Curves.eps}
\caption{ROC Curves for Ensemble Models.}
\end{figure}
\begin{figure}[H]
\centering
\includegraphics[width=0.8\linewidth]{Ex6/Confusion_Matrices.eps}
\caption{Confusion Matrices comparing errors.}
\end{figure}

\section{Discussion}
Bagging effectively reduced the variance of the base Decision Tree. Boosting models (AdaBoost and Gradient Boosting) improved upon bias by sequentially focusing on misclassified samples. The Stacked Ensemble leveraged diverse base estimators (SVM, Naive Bayes, Decision Tree) and optimally combined their predictive power using a Logistic Regression meta-learner.

\section{Conclusion}
Ensemble techniques consistently outperformed standalone estimators. Gradient Boosting and Stacking demonstrated superior generalization on the test set, exhibiting the lowest false negative rates, which is critical for cancer diagnosis. The experiment practically illustrated the theoretical trade-offs between variance reduction (Bagging) and bias reduction (Boosting).

\section{References}
Use IEEE style.

\end{document}
